\chapter{Implementation and Results}
\begin{spacing}{1.2}
\minitoc
\thispagestyle{MyStyle}
\end{spacing}
\newpage

\section{Introduction}

This chapter presents the complete implementation of the Time Guide AI system and the results obtained through its development and deployment. It describes the three-tier architecture, database design, backend services, frontend user interface, artificial intelligence integration, data import pipeline, security mechanisms, and testing strategies. Additionally, it discusses the user experience design, including the conversational interface, sidebar navigation, suggestion system, and captured demonstrations from the deployed chatbot.

The implementation demonstrates how theoretical foundations presented in the literature chapter are translated into a coherent, production-ready academic information system that balances technical sophistication with practical usability.

\section{Global System Architecture}

Time Guide AI implements a classical three-tier architecture, designed for modularity, scalability, and clear separation of concerns. This architecture ensures that frontend, backend, and database components can evolve independently while maintaining consistent interfaces.

\begin{figure}[H]
    \centering
    \includegraphics[width=1\linewidth]{images/architecture_diagram.png}
    \caption{Three-tier system architecture showing the presentation layer (React frontend), application layer (FastAPI backend with modular services), and data layer (PostgreSQL database). Communication flows through REST API endpoints.}
    \label{fig:global-architecture}
\end{figure}

\subsection{Presentation Tier: React Frontend}

The presentation tier is implemented using React 18 with TypeScript, providing a modern, responsive web interface accessible from any browser. The frontend is built with the following technologies:

\begin{itemize}
    \item \textbf{React 18}: Component-based UI framework enabling interactive user interfaces.
    \item \textbf{TypeScript}: Provides static type checking and improved code maintainability.
    \item \textbf{Vite}: Modern build tool offering fast development server startup and optimized production builds.
    \item \textbf{Tailwind CSS}: Utility-first CSS framework for rapid, responsive UI development.
    \item \textbf{shadcn/ui}: Pre-built, accessible React components (buttons, dialogs, forms, cards) styled with Tailwind CSS.
    \item \textbf{Framer Motion}: Animation library for smooth transitions and visual feedback.
    \item \textbf{React Router DOM}: Client-side routing for multi-page navigation.
    \item \textbf{TanStack Query}: Server state management for efficient API data fetching and caching.
\end{itemize}

The frontend consists of five main pages:

\begin{enumerate}
    \item \textbf{Index Page}: Landing page with system overview and call-to-action buttons.
    \item \textbf{Chat Page}: Main user interface for conversational interaction with the system.
    \item \textbf{Auth Page}: Authentication interface for login and signup.
    \item \textbf{Admin Page}: Administration dashboard for timetable import and management.
    \item \textbf{NotFound Page}: 404 error page for undefined routes.
\end{enumerate}

\subsection{Application Tier: FastAPI Backend}

The application tier is implemented using FastAPI, a modern Python web framework designed for building high-performance REST APIs. FastAPI provides several advantages:

\begin{itemize}
    \item \textbf{Asynchronous Processing}: Native support for async/await, enabling concurrent request handling.
    \item \textbf{Automatic Documentation}: FastAPI generates OpenAPI/Swagger documentation automatically.
    \item \textbf{Type Safety}: Full support for Python type hints with automatic request validation.
    \item \textbf{Dependency Injection}: Clean, declarative dependency management through the Depends() mechanism.
\end{itemize}

The backend is organized into three logical layers:

\begin{enumerate}
    \item \textbf{Routes Layer} (\texttt{app/routes/}): Handles HTTP requests and exposes API endpoints for authentication, chat, and administration.
    \item \textbf{Services Layer} (\texttt{app/services/}): Implements core business logic through modular, specialized services.
    \item \textbf{Models Layer} (\texttt{app/models/}): Defines database models using SQLAlchemy ORM.
\end{enumerate}

API endpoints are organized by functionality:

\begin{itemize}
    \item \textbf{Authentication Endpoints}:
    \begin{itemize}
        \item \texttt{POST /api/auth/login}: User login with email and password.
        \item \texttt{POST /api/auth/signup}: User registration.
        \item \texttt{GET /api/auth/me}: Retrieve current authenticated user info.
        \item \texttt{POST /api/auth/logout}: Terminate user session.
    \end{itemize}
    
    \item \textbf{Chat Endpoint}:
    \begin{itemize}
        \item \texttt{POST /api/chat}: Process user message and return AI-generated response.
    \end{itemize}
    
    \item \textbf{Admin Endpoints}:
    \begin{itemize}
        \item \texttt{POST /api/admin/imports/upload}: Upload timetable Excel files for import.
        \item \texttt{GET /api/admin/imports/status}: Retrieve import history and status.
    \end{itemize}
    
    \item \textbf{Health Check}:
    \begin{itemize}
        \item \texttt{GET /health}: System health verification endpoint.
    \end{itemize}
\end{itemize}

\subsection{Data Tier: PostgreSQL Database}

The data tier uses PostgreSQL 14+, a robust relational database management system. PostgreSQL provides:

\begin{itemize}
    \item \textbf{ACID Compliance}: Strong consistency guarantees for transactional integrity.
    \item \textbf{Advanced Indexing}: Support for btree, hash, and other index types for query optimization.
    \item \textbf{Foreign Key Constraints}: Referential integrity enforcement across related tables.
    \item \textbf{JSON Support}: Native handling of semi-structured data when needed.
\end{itemize}

\section{Database Design}

The database schema models the academic structure of ENET'Com, including students, teachers, courses, rooms, timetables, and calendar events. The design follows normalization principles to minimize redundancy and ensure consistency.

\begin{figure}[H]
    \centering
    \includegraphics[width=1\linewidth]{images/database_schema.png}
    \caption{PostgreSQL database schema showing relationships between academic entities. Core tables include classes, professors, rooms, sessions, and versions. Authentication tables manage user access control.}
    \label{fig:database-schema}
\end{figure}

\subsection{Academic Entity Tables}

\begin{center}
\begin{tabular}{|p{2.5cm}|p{10cm}|}
\hline
\textbf{Table} & \textbf{Purpose and Structure} \\
\hline

\textbf{annees\_universitaires} & 
Stores academic years (e.g., 2025-2026). Fields: id (PK), libelle (year label), date\_debut (start date), date\_fin (end date). \\
\hline

\textbf{semestres} & 
Represents semesters (S1, S2). Fields: id (PK), nom (name), annee\_id (FK to academic year). Each semester belongs to an academic year. \\
\hline

\textbf{periodes} & 
Defines periods within semesters (P1, P2). Fields: id (PK), nom (period name), semestre\_id (FK), date\_debut, date\_fin. Periods mark distinct teaching phases. \\
\hline

\textbf{departements} & 
Academic departments (Informatique, Électronique, etc.). Fields: id (PK), nom (department name). \\
\hline

\textbf{classes} & 
Student classes (e.g., 2 ING GII 3, 1 TIC 1). Fields: id (PK), nom (class name), departement\_id (FK), semestre\_id (FK). \\
\hline

\textbf{professeurs} & 
Teaching staff. Fields: id (PK), nom\_complet (full name), grade (professor rank), specialite (specialization). Indexed on nom\_complet for efficient lookup. \\
\hline

\textbf{matieres} & 
Course subjects (e.g., TRAIT IMAGES, BASES DE DONNÉES). Fields: id (PK), nom (subject name), code (course code). \\
\hline

\textbf{salles} & 
Classrooms and teaching rooms. Fields: id (PK), nom (room name/code like C14), type (room type), capacite (capacity). \\
\hline

\textbf{groupes} & 
Student groups within classes. Fields: id (PK), nom (group name), classe\_id (FK). Enables sub-class scheduling. \\
\hline

\textbf{emplois\_versions} & 
Timetable version management. Fields: id (PK), classe\_id (FK), version\_date (import date), actif (boolean). Critical for maintaining multiple versions and ensuring only active timetables are queried. \\
\hline

\textbf{seances} & 
Individual teaching sessions. Fields: id (PK), version\_id (FK), classe\_id (FK), matiere\_id (FK), professeur\_id (FK), salle\_id (FK), groupe\_id (FK, nullable), periode\_id (FK, nullable), jour (day name), heure\_debut (start time), heure\_fin (end time), type\_seance (cours/TD/TP). Indexed on jour, classe\_id, version\_id for query performance. \\
\hline

\textbf{emplois\_enseignants\_seances} & 
Denormalized cache of teacher schedules for efficient queries. Fields: id (PK), semestre\_id (FK), professeur\_nom\_complet, classe\_nom, matiere\_nom, salle\_nom, jour, heure\_debut, heure\_fin, periode\_nom, type\_seance. This table enables fast teacher-specific queries without complex joins. \\
\hline

\textbf{vacances\_jours\_feries} & 
Academic calendar events. Fields: id (PK), nom (event name), date\_debut, date\_fin, type (vacances/jour\_ferie/examen/revision), annee\_id (FK). Supports holiday and exam period awareness. \\
\hline

\end{tabular}
\end{center}

\subsection{Authentication and Security Tables}

\begin{center}
\begin{tabular}{|p{2.5cm}|p{10cm}|}
\hline
\textbf{Table} & \textbf{Purpose and Structure} \\
\hline

\textbf{auth\_users} & 
User account storage. Fields: id (PK, unique token), email (unique, indexed), full\_name, password\_hash (PBKDF2-HMAC-SHA256), role (student/professor/admin), created\_at. Passwords are never stored in plaintext. \\
\hline

\textbf{auth\_sessions} & 
Active user sessions. Fields: id (PK), user\_id (FK), email, full\_name, role, token\_hash (SHA256 of JWT), expires\_at, created\_at. Sessions expire after 14 days by default (configurable via AUTH\_SESSION\_TTL\_DAYS). \\
\hline

\end{tabular}
\end{center}

\subsection{Database Consistency and Integrity}

The schema enforces referential integrity through cascade relationships:

\begin{itemize}
    \item A test session (seance) references a class, professor, room, and subject. All must exist.
    \item A class belongs to a semester, which belongs to an academic year.
    \item Only active timetable versions are used for query responses, preventing stale data.
    \item Unique indexes on email addresses prevent duplicate user accounts.
\end{itemize}

\section{Backend Services Architecture}

The backend is organized into modular services, each responsible for a specific domain. This service-oriented architecture enables independent testing, maintenance, and evolution of components.

\begin{figure}[H]
    \centering
    \includegraphics[width=1\linewidth]{images/services_architecture.png}
    \caption{Backend services architecture showing the modular organization of core services. Each service handles a specific responsibility: routing, SQL generation, authentication, file import, and external data retrieval.}
    \label{fig:services-architecture}
\end{figure}

\subsection{Core Services}

\subsubsection{Intent Router Service}
The Intent Router is the entry point for query processing. It classifies user intent, extracts relevant entities, and decides the execution path. The classification hierarchy is:

\begin{enumerate}
    \item \textbf{Conversation Markers}: Detect greetings, thanks, help requests using keyword matching.
    \item \textbf{Deterministic Classification}: Use academic keywords and patterns to classify intent with high confidence.
    \item \textbf{History-Based Reasoning}: Consider conversation history for context-aware follow-ups.
    \item \textbf{Groq LLM Classification}: Invoke language model when deterministic methods lack confidence.
    \item \textbf{Fallback}: Request clarification if classification confidence is below threshold (0.75).
\end{enumerate}

Intent categories recognized are:

\begin{itemize}
    \item CLASS\_SCHEDULE, PROF\_SCHEDULE, PROF\_LOCATION, PROF\_CLASS
    \item ROOM\_SCHEDULE, ROOM\_AVAILABLE, AVAILABLE\_ROOMS
    \item CALENDAR, UNIVERSITY\_INFO
    \item CHAT\_SMALLTALK, OUT\_OF\_SCOPE, NEEDS\_CLASS, NEEDS\_PROF\_CONFIRMATION
\end{itemize}

\subsubsection{SQL Agent Service}
The SQL Agent handles all database query generation and execution for timetable questions. It operates in phases:

\begin{enumerate}
    \item \textbf{Academic Context Retrieval}: Query the database for the current academic year, semester, and period.
    \item \textbf{Entity Validation}: Verify that extracted entities (professor names, room codes, class names) exist in the database.
    \item \textbf{Missing Information Check}: If required information is unavailable (e.g., professor name not specified for a teacher query), request clarification.
    \item \textbf{SQL Generation}: Generate or select a parameterized SQL query based on intent and extracted entities.
    \item \textbf{Query Validation}: Ensure query is SELECT-only, references valid tables, and uses valid column names.
    \item \textbf{Execution}: Execute query on PostgreSQL with parameters to prevent SQL injection.
    \item \textbf{Result Formatting}: Transform query results into natural language responses.
\end{enumerate}

Schema information is provided to the Language Model before SQL generation to prevent hallucination:

\begin{verbatim}
Timetables are queried from seances table.
Active timetable versions are identified via:
  JOIN emplois_versions v ON v.id = s.version_id
  WHERE v.actif = true AND v.classe_id = s.classe_id

Common filters:
- By professor: WHERE p.nom_complet LIKE %name%
- By room: WHERE s.salle_id = %room_id%
- By day: WHERE s.jour = %day%
- By time: WHERE s.heure_debut <= %time% AND s.heure_fin > %time%
\end{verbatim}

\subsubsection{Groq Service}
The Groq Service provides language model capabilities for tasks beyond deterministic logic:

\begin{itemize}
    \item \textbf{Intent Classification}: When deterministic classification is uncertain.
    \item \textbf{Entity Extraction}: Identifying names, times, dates from unstructured input.
    \item \textbf{Response Formatting}: Converting raw database results to natural language.
    \item \textbf{Smalltalk Generation}: Producing conversational responses for greetings and help requests.
    \item \textbf{Clarification Questions}: Asking natural-language disambiguation questions.
    \item \textbf{Fallback SQL Generation}: Generating custom SQL for query patterns not pre-defined.
\end{itemize}

The service uses \texttt{llama-3.3-70b-versatile} with temperature 0.0 for deterministic tasks (classification) and temperature 0.1-0.3 for varied but sensible responses (formatting, smalltalk).

\subsubsection{University Information Service}
This service provides institutional information grounded in official ENET'Com web pages:

\begin{itemize}
    \item Fetches information from the ENET'Com website for queries about departments, programs, contact information, research laboratories, internship opportunities.
    \item Caches responses (TTL: 30 minutes) to reduce external requests.
    \item Performs web scraping and HTML parsing to extract structured academic data.
    \item Maps keywords to relevant website pages for targeted retrieval.
\end{itemize}

\subsubsection{Auth Service}
Manages user authentication and authorization:

\begin{itemize}
    \item \textbf{Password Hashing}: Uses PBKDF2-HMAC-SHA256 with 200,000 iterations and a 16-byte salt.
    \item \textbf{Login/Signup}: Creates user accounts and validates credentials.
    \item \textbf{JWT Token Generation}: Issues digitally signed tokens with user ID, role, and expiration.
    \item \textbf{Session Management}: Tracks active sessions and enforces role-based access control.
    \item \textbf{Token Verification}: Validates JWT signatures and checks expiration before processing requests.
\end{itemize}

\subsubsection{Admin Import Service}
Orchestrates the complete timetable import workflow:

\begin{itemize}
    \item \textbf{File Type Detection}: Automatically identifies workbook category (student\_s1, teachers\_s2, rooms\_s1, calendar, etc.).
    \item \textbf{Validation}: Checks semester consistency and category correctness.
    \item \textbf{Parsing}: Extracts timetable data from Excel using the Excel Parser.
    \item \textbf{Transformation}: Normalizes extracted data (resolves professor names, room codes, subject names to database IDs).
    \item \textbf{Loading}: Inserts records into the database within a transaction.
    \item \textbf{Version Management}: Marks new timetable as active and old versions as inactive.
    \item \textbf{Status Reporting}: Returns detailed feedback including session count, errors, and warnings.
\end{itemize}

\subsubsection{Excel Parser Service}
Converts Excel workbooks into structured timetable records:

\begin{itemize}
    \item \textbf{Format Recognition}: Identifies row/column structure (time slots, days).
    \item \textbf{Cell Block Parsing}: Extracts multi-line cell content (subject, professor, room, group, period marker).
    \item \textbf{Time Extraction}: Uses regex to parse time ranges (e.g., ``08:00 - 10:00'').
    \item \textbf{Period Detection}: Recognizes (P1) and (P2) markers in cell content.
    \item \textbf{Type Inference}: Classifies sessions as cours (lecture), TD (tutorial), or TP (practical).
    \item \textbf{Entity Resolution}: Matches extracted professor names to database records using fuzzy matching (Levenshtein distance).
\end{itemize}

\section{Frontend User Interface Design}

The frontend provides an intuitive, modern interface for users to interact with the system. The design emphasizes clarity, responsiveness, and accessibility.

\subsection{Page Structure}

\subsubsection{Chat Page}

The Chat page is the main user-facing interface. It consists of several key components:

\begin{figure}[H]
    \centering
    \includegraphics[width=1\linewidth]{images/chat_interface_layout.png}
    \caption{Chat interface layout showing sidebar, main conversation area, welcome screen with suggestions, and chat input. Sidebar displays conversation history and allows quick access to previous discussions.}
    \label{fig:chat-interface}
\end{figure}

\paragraph{Header Section}
\begin{itemize}
    \item Application branding (Enet'Bot logo and title).
    \item Subtitle: ``Conversation universitaire ENET'Com''.
    \item Status indicator: ``Posez vos questions sur les classes, salles, enseignants et calendrier'' (Ask questions about classes, rooms, teachers, and calendar).
    \item Admin access button (for administrators only).
    \item Logout button.
\end{itemize}

\paragraph{Sidebar (Conversation History)}
The sidebar displays previous conversations and provides navigation:

\begin{itemize}
    \item \textbf{New Chat Button}: Creates a new conversation session.
    \item \textbf{Conversation List}: Shows up to 10 most recent conversations, each with:
    \begin{itemize}
        \item Conversation title (first 30 characters of the first user message).
        \item Conversation date.
        \item Delete button for individual conversations.
    \end{itemize}
    \item \textbf{Local Storage Persistence}: Conversations are saved to browser localStorage, keyed by user ID.
    \item \textbf{Responsive Behavior}: On mobile devices (width \textless 1024px), sidebar toggles on/off; on desktop, it's always visible.
\end{itemize}

Implementation details:

\begin{verbatim}
// ChatSidebar.tsx: Manages conversation history
- Saves conversations to localStorage: localStorage.setItem(storageKey, JSON.stringify(conversations))
- Loads on component mount: const saved = localStorage.getItem(storageKey)
- Limits to 10 most recent conversations: const sliced = next.slice(0, 10)
- Auto-saves conversation updates in real-time as messages are added
\end{verbatim}

\paragraph{Welcome Screen (First Load)}
When a conversation is empty, the Welcome Screen displays:

\begin{itemize}
    \item System greeting and description.
    \item \textbf{Suggestion Buttons}: Quick-access buttons for commonly asked question types.
    \item \textbf{Example Questions}: Pre-formatted suggestions guiding users on how to interact.
\end{itemize}

Example suggestions displayed:

\begin{enumerate}
    \item ``Quel est mon emploi du temps demain?'' (What is my timetable tomorrow?)
    \item ``Où est Mr BEN SLIMA maintenant?'' (Where is Mr BEN SLIMA now?)
    \item ``Quelles salles sont libres lundi?'' (Which rooms are free Monday?)
    \item ``Quand commencent les vacances?'' (When do holidays start?)
    \item ``Quelles sont les coordonnées du département informatique?'' (What is the contact info for the IT department?)
\end{enumerate}

These suggestions guide users on the types of queries the system supports and encourage natural-language interaction patterns.

\paragraph{Conversation Area}
The main conversation area displays messages in a scrollable list:

\begin{itemize}
    \item \textbf{User Messages}: Displayed on the right side with user styling, timestamp.
    \item \textbf{Assistant Messages}: Displayed on the left side with bot styling, timestamp.
    \item \textbf{Typing Indicator}: Animated dots while the system processes a query.
    \item \textbf{Auto-scroll}: Automatically scrolls to the latest message as new messages arrive.
\end{itemize}

\paragraph{Chat Input Component}
At the bottom, the input field allows message submission:

\begin{itemize}
    \item Text input with placeholder: ``Posez votre question...''.
    \item Send button (arrow icon) triggered on button click or Enter key press.
    \item Character limit feedback (optional).
    \item Disabled state while processing (isLoading flag prevents duplicate submissions).
\end{itemize}

\subsection{Message Request-Response Flow}

Users submit messages through the following sequence:

\begin{enumerate}
    \item User types a message in the input field.
    \item User presses Enter or clicks Send button.
    \item Frontend captures the message and extracts conversation history (last 3 messages for context).
    \item Frontend performs a POST request to \texttt{/api/chat}:
    \begin{verbatim}
    POST /api/chat
    {
      "message": "Quel est l'emploi de temps de 2 GII 3?",
      "user_role": "student",
      "user_class": null,
      "history": [
        {"role": "user", "content": "Bonjour"},
        {"role": "assistant", "content": "Bonjour!"}
      ]
    }
    \end{verbatim}
    \item Backend processes the request (see Section \ref{sec:backend-processing}).
    \item Backend returns ChatResponse with the generated answer:
    \begin{verbatim}
    {
      "response": "Lundi: TRAIT IMAGES 08:00-10:00..."
    }
    \end{verbatim}
    \item Frontend displays the response in the conversation area.
    \item Frontend saves the conversation to localStorage for persistence.
\end{enumerate}

\subsection{Admin Dashboard Page}

The Admin page provides tools for managing timetable imports:

\begin{figure}[H]
    \centering
    \includegraphics[width=1\linewidth]{images/admin_dashboard.png}
    \caption{Admin dashboard showing upload cards for each workbook category, latest file information, and import status indicators with success/error notifications.}
    \label{fig:admin-dashboard}
\end{figure}

\paragraph{Upload Categories}
The dashboard displays cards for each upload category:

\begin{center}
\begin{tabular}{|p{3cm}|p{2cm}|p{8cm}|}
\hline
\textbf{Category} & \textbf{Icon} & \textbf{Description} \\
\hline
student\_s1 & GraduationCap & Student timetables for Semester 1 \\
student\_s2 & GraduationCap & Student timetables for Semester 2 \\
teachers\_s1 & Users & Teacher timetables for Semester 1 \\
teachers\_s2 & Users & Teacher timetables for Semester 2 \\
rooms\_s1 & DoorOpen & Room schedules for Semester 1 \\
rooms\_s2 & DoorOpen & Room schedules for Semester 2 \\
calendar & CalendarDays & Academic calendar (holidays, exams, etc.) \\
\hline
\end{tabular}
\end{center}

\paragraph{Card Features}
Each upload card includes:

\begin{itemize}
    \item Category icon and title.
    \item Current latest file information (if any):
    \begin{itemize}
        \item Filename.
        \item Upload timestamp.
        \item File size.
    \end{itemize}
    \item \textbf{File Input}: Click to select and upload an Excel file.
    \item \textbf{Status Indicator}: Visual feedback (success/error/pending).
    \item \textbf{Notification}: Toast notification displayed after upload (success or error message).
\end{itemize}

\paragraph{Import Status Display}
After each upload, the dashboard displays:

\begin{itemize}
    \item Category name.
    \item Upload status (success/error).
    \item Number of sessions imported (for successful uploads).
    \item Error message or warning (if applicable).
\end{itemize}

\section{End-to-End Query Processing Workflow}

\subsection{Query Submission to Response Generation}

The complete workflow from user query to system response is as follows:

\begin{figure}[H]
    \centering
    \includegraphics[width=1\linewidth]{images/query_processing_flow.png}
    \caption{End-to-end workflow showing the transformation of a user query through intent classification, entity extraction, execution path selection, and response generation.}
    \label{fig:query-processing-flow}
\end{figure}

\begin{enumerate}
    \item \textbf{User Input Submission}: User submits a natural-language question from the Chat interface.
    
    \item \textbf{Frontend Packaging}: Chat.tsx packages the message, user role, class, and conversation history into a ChatRequest.
    
    \item \textbf{API Request}: Frontend sends a POST request to \texttt{/api/chat} endpoint.
    
    \item \textbf{Backend Reception}: FastAPI route handler receives the request and creates a database session.
    
    \item \textbf{Intent Classification}: IntentRouter.route() is invoked:
    \begin{itemize}
        \item Normalize message (remove extra spaces, convert to lowercase).
        \item Check for conversation markers (greetings, thanks) → if matched, route to SMALLTALK.
        \item Check for academic markers (emploi, prof, salle, etc.) → classify as ACADEMIC or NON\_ACADEMIC.
        \item Extract entities (professor names, room codes, class names, day/time indicators).
        \item Determine intent category (CLASS\_SCHEDULE, PROF\_LOCATION, ROOM\_AVAILABLE, etc.).
        \item Calculate confidence score.
    \end{itemize}
    
    \item \textbf{Execution Path Selection}: Based on confidence and intent, select execution target:
    \begin{itemize}
        \item \textbf{SQL\_AGENT}: For timetable questions (database queries).
        \item \textbf{UNIVERSITY\_SERVICE}: For institutional information requests.
        \item \textbf{SMALLTALK}: For conversational messages.
        \item \textbf{CLARIFICATION}: If disambiguation is needed.
        \item \textbf{OUT\_OF\_SCOPE}: If query is unrelated to system capabilities.
    \end{itemize}
    
    \item \textbf{Timetable Query Execution} (if SQL\_AGENT selected):
    \begin{itemize}
        \item SQL Agent validates extracted entities against the database.
        \item If required information is missing, a clarification message is returned.
        \item Otherwise, SQL Agent generates or selects a parameterized SQL query.
        \item Query is validated to ensure it's SELECT-only.
        \item Query is executed on PostgreSQL.
        \item Results are formatted into a natural-language response.
    \end{itemize}
    
    \item \textbf{Institutional Query Execution} (if UNIVERSITY\_SERVICE selected):
    \begin{itemize}
        \item University Information Service maps the query to relevant website pages.
        \item Fetches content from ENET'Com website (with caching).
        \item Extracts relevant information and formats response.
    \end{itemize}
    
    \item \textbf{Response Generation}:
    \begin{itemize}
        \item Format results according to query type (timetable, text, list, etc.).
        \item If needed, polish response using Groq for natural language formatting.
        \item Return ChatResponse object.
    \end{itemize}
    
    \item \textbf{Response Display}: Frontend receives the response and displays it in the Chat interface.
    
    \item \textbf{Conversation Persistence}: Frontend saves the updated conversation (including both user message and AI response) to localStorage.
\end{enumerate}

\section{Data Import and ETL Pipeline}

The Excel import pipeline enables administrators to keep the system's knowledge base synchronized with operational timetable data. The process is designed to minimize errors and maximize data quality.

\subsection{Import Workflow}

\begin{fig}[H]
    \centering
    \includegraphics[width=1\linewidth]{images/import_workflow.png}
    \caption{ETL pipeline for Excel timetable import. Shows file upload, category detection, validation, parsing, transformation, and database loading phases.}
    \label{fig:import-workflow}
\end{figure}

\paragraph{Phase 1: File Reception}
\begin{itemize}
    \item Administrator selects an Excel file using the Admin dashboard.
    \item File is uploaded to the server and temporarily stored.
    \item File size and format are validated (must be .xlsx).
\end{itemize}

\paragraph{Phase 2: Category Detection}
The admin\_import\_service examines the file structure to determine its category:

\begin{itemize}
    \item Check sheet names for patterns indicating workbook type.
    \item Examine column headers for day names (Lundi, Mardi, etc.).
    \item Look for semester markers (S1, S2, Period tags like P1, P2).
    \item Analyze content patterns (room vs. class vs. teacher naming conventions).
    \item Classify as: student\_s1, student\_s2, teachers\_s1, teachers\_s2, rooms\_s1, rooms\_s2, or calendar.
\end{itemize}

\paragraph{Phase 3: Validation}
Before parsing, the system validates consistency:

\begin{itemize}
    \item \textbf{Category Match}: Verify that the detected category matches the expected type.
    \item \textbf{Semester Consistency}: If uploading student\_s1, all classes must belong to Semester 1.
    \item \textbf{Entity Existence}: Verify that referenced professors, rooms, and subjects exist in the database.
    \item \textbf{Date Range}: Verify that import date falls within a valid academic year.
\end{itemize}

\paragraph{Phase 4: Parsing and Transformation}
The Excel Parser extracts and normalizes data:

\begin{enumerate}
    \item Extract time slots from row headers using regex: \texttt{(\textbackslash d\{1,2\})[h:].*?-.*?\textbackslash d\{1,2\}}.
    \item Extract day names from column headers.
    \item For each cell:
    \begin{itemize}
        \item Split multi-line content by newlines.
        \item Identify subject name (usually first line).
        \item Extract period marker (P1) or (P2) if present.
        \item Identify professor line (preceded by Mr, Mme, Dr, Prof).
        \item Identify room line (format like C14, A05).
        \item Identify group/audience (e.g., Master students, Groupe 1).
    \end{itemize}
    \item Infer session type (cours, TD, TP) from subject name or context.
\end{enumerate}

\paragraph{Phase 5: Entity Resolution}
Extracted entities are linked to database records:

\begin{itemize}
    \item \textbf{Professor Resolution}: Match extracted name to \texttt{professeurs} table using:
    \begin{itemize}
        \item Exact match.
        \item Partial/fuzzy match (Levenshtein distance for typo tolerance).
        \item If no match, create new professor record or skip session with warning.
    \end{itemize}
    
    \item \textbf{Room Resolution}: Match room code to \texttt{salles} table.
    \item \textbf{Subject Resolution}: Match subject name to \texttt{matieres} table.
    \item \textbf{Class Resolution}: Match class name to \texttt{classes} table.
    \item \textbf{Group Resolution}: Create group records if needed.
\end{itemize}

\paragraph{Phase 6: Database Loading}
Records are inserted within a single transaction to ensure atomicity:

\begin{verbatim}
Transaction Process:
1. BEGIN TRANSACTION
2. FOR each parsed session:
   a. Check for exact duplicate (prevent re-import of identical session)
   b. IF new session:
      - INSERT into seances table
   c. ELSE:
      - SKIP (duplicate)
3. UPDATE emplois_versions SET actif = true WHERE id = new_version
4. UPDATE emplois_versions SET actif = false WHERE classe_id = class 
                                                  AND id != new_version
5. IF all inserts succeed:
   - COMMIT transaction
   - RETURN success with session count
6. IF any error:
   - ROLLBACK all changes
   - RETURN error message
\end{verbatim}

\paragraph{Phase 7: Status Reporting}
The admin receives detailed feedback:

\begin{verbatim}
{
  "success": true,
  "category": "student_s1",
  "sessions_imported": 48,
  "sessions_updated": 2,
  "sessions_skipped": 0,
  "version_id": 15,
  "import_date": "2026-02-15"
}
\end{verbatim}

\section{Security and Authentication}

The system implements practical security controls appropriate for an academic information system handling operational data.

\subsection{Authentication Mechanism}

\subsubsection{Password Management}
\begin{itemize}
    \item \textbf{Hashing Algorithm}: PBKDF2-HMAC-SHA256 with 200,000 iterations.
    \item \textbf{Salt Generation}: 16-byte random salt generated per password using \texttt{secrets.token\_hex(16)}.
    \item \textbf{Storage Format}: \texttt{\$salt\$hash} (salt and hash separated by \$).
    \item \textbf{Comparison}: Uses \texttt{secrets.compare\_digest()} for constant-time comparison to resist timing attacks.
\end{itemize}

\subsubsection{JWT Tokens}
\begin{itemize}
    \item \textbf{Token Generation}: Digitally signed JWT containing user ID, email, role, and expiration timestamp.
    \item \textbf{Expiration}: 14 days (configurable via AUTH\_SESSION\_TTL\_DAYS environment variable).
    \item \textbf{Session Tracking}: Token hash stored in \texttt{auth\_sessions} table to prevent token reuse after logout.
    \item \textbf{Verification}: Every API request validates JWT signature and checks expiration.
\end{itemize}

\subsection{Authorization and Access Control}

User roles and permissions:

\begin{center}
\begin{tabular}{|p{2.5cm}|p{11cm}|}
\hline
\textbf{Role} & \textbf{Permissions} \\
\hline
\textbf{student} & 
Can access chat interface, ask timetable questions, view information. Cannot access admin functions or modify data. \\
\hline
\textbf{professor} & 
Can access chat interface, query their own teaching schedule and classes. May have limited access to student information. \\
\hline
\textbf{admin} & 
Can access admin dashboard, upload Excel workbooks, detect categories, manage timetable versions, view import status. Can override data or reassign classes if needed. \\
\hline
\end{tabular}
\end{center}

Access control is enforced at the route level in FastAPI using the \texttt{Depends()} mechanism:

\begin{verbatim}
@router.post("/api/admin/imports/upload")
async def admin_import_upload(
    db: Session = Depends(get_db),
    _admin_user=Depends(get_admin_user),  # Enforces admin role
):
    # Only admin users can reach this function
    ...

def get_admin_user(current_user=Depends(get_current_user)):
    if current_user.role != "admin":
        raise HTTPException(status_code=403, detail="Admin access required")
    return current_user
\end{verbatim}

\subsection{SQL Injection Prevention}

All database queries use parameterized statements:

\begin{verbatim}
# Safe (parameterized):
result = db.execute(
    text("SELECT * FROM seances WHERE professeur_id = :prof_id"),
    {"prof_id": prof_id}
)

# Dangerous (string interpolation - NEVER used):
query = f"SELECT * FROM seances WHERE professeur_id = {prof_id}"  # WRONG!
\end{verbatim}

\subsection{Read-Only Enforcement}

Generated SQL queries are validated before execution:

\begin{itemize}
    \item Regex check: Query must match \texttt{\^\\s*SELECT\\s+} and must NOT contain INSERT, UPDATE, DELETE, or DROP.
    \item Parser verification: A lightweight SQL parser confirms SELECT is the first clause.
    \item Database permission: Non-admin database users lack INSERT, UPDATE, DELETE privileges.
\end{itemize}

\section{API Design and Specification}

The API follows RESTful principles with JSON request/response payloads. All endpoints use CORS-enabled communication between frontend and backend.

\subsection{Authentication Endpoints}

\begin{center}
\begin{tabular}{|p{2.5cm}|p{2.5cm}|p{8cm}|}
\hline
\textbf{Endpoint} & \textbf{Method} & \textbf{Purpose} \\
\hline
\texttt{/api/auth/login} & POST & 
Authenticate user. Request: \{email, password\}. Response: \{token, user\}. \\
\hline
\texttt{/api/auth/signup} & POST & 
Register new user. Request: \{email, password, full\_name\}. Response: \{user\}. \\
\hline
\texttt{/api/auth/me} & GET & 
Retrieve current authenticated user. Authorization required. Response: \{user\}. \\
\hline
\texttt{/api/auth/logout} & POST & 
Terminate session. Authorization required. Response: \{success: true\}. \\
\hline
\end{tabular}
\end{center}

\subsection{Chat Endpoint}

\begin{verbatim}
POST /api/chat

Request:
{
  "message": string,              // User's natural language question
  "user_role": string,            // "student", "professor", "admin"
  "user_class": string | null,    // User's class (optional)
  "history": Message[]            // Previous messages [{"role": "user"/"assistant", "content": "..."}]
}

Response:
{
  "response": string              // System's generated response
}
\end{verbatim}

\subsection{Admin Endpoints}

\begin{verbatim}
POST /api/admin/imports/upload

Request (multipart/form-data):
- student_s1: File | null
- student_s2: File | null
- rooms_s1: File | null
- rooms_s2: File | null
- teachers_s1: File | null
- teachers_s2: File | null
- calendar: File | null

Response:
{
  "results": [
    {
      "category": "student_s1",
      "status": "success" | "error",
      "message": string,
      "sessions_imported": number | null
    },
    ...
  ]
}

---

GET /api/admin/imports/status

Response:
{
  "calendar_warning": string | null,
  "categories": [
    {
      "key": "student_s1",
      "label": "Etudiants S1",
      "latest_file": {
        "filename": string,
        "uploaded_at": datetime,
        "size_label": string
      } | null
    },
    ...
  ]
}
\end{verbatim}

\section{Results and Demonstration}

The Time Guide AI system has been successfully implemented and deployed, demonstrating all designed functionality. This section presents captured demonstrations of the system in action.

\subsection{Chat Interface Examples}

The following screenshots capture typical user interactions with the system:

\begin{figure}[H]
    \centering
    \includegraphics[width=1\linewidth]{images/chat_demo_01_welcome.png}
    \caption{Welcome screen showing the chat interface with suggested questions. Users can click suggested questions or type their own queries in natural language.}
    \label{fig:chat-demo-welcome}
\end{figure}

\begin{figure}[H]
    \centering
    \includegraphics[width=1\linewidth]{images/chat_demo_02_class_schedule.png}
    \caption{Example query: ``Quel est l'emploi de temps de 2 ING GII 3?" (What is the timetable for class 2 ING GII 3?). System responses with detailed schedule showing times, professors, rooms, and subjects.}
    \label{fig:chat-demo-class-schedule}
\end{figure}

\begin{figure}[H]
    \centering
    \includegraphics[width=1\linewidth]{images/chat_demo_03_professor_location.png}
    \caption{Example query: ``Où est Mr BEN SLIMA maintenant?" (Where is Mr BEN SLIMA now?). System retrieves current or upcoming professor location in real-time.}
    \label{fig:chat-demo-professor-location}
\end{figure}

\begin{figure}[H]
    \centering
    \includegraphics[width=1\linewidth]{images/chat_demo_04_room_availability.png}
    \caption{Example query: ``Quelles salles sont libres lundi?" (Which rooms are free Monday?). System lists available rooms with time windows.}
    \label{fig:chat-demo-room-availability}
\end{figure}

\begin{figure}[H]
    \centering
    \includegraphics[width=1\linewidth]{images/chat_demo_05_institutional_info.png}
    \caption{Example query: ``Quelles sont les coordonnées du département informatique?" (What is the contact info for the IT department?). System provides grounded information from ENET'Com website.}
    \label{fig:chat-demo-institutional}
\end{figure}

\begin{figure}[H]
    \centering
    \includegraphics[width=1\linewidth]{images/chat_demo_06_conversation_history.png}
    \caption{Sidebar showing conversation history with multiple past discussions. Users can switch between conversations or start a new chat.}
    \label{fig:chat-demo-history}
\end{figure}

\subsection{Admin Dashboard Examples}

\begin{figure}[H]
    \centering
    \includegraphics[width=1\linewidth]{images/admin_demo_01_upload_interface.png}
    \caption{Admin dashboard showing upload cards for each workbook category. Administrators can drag-and-drop or click to select Excel files for import.}
    \label{fig:admin-demo-upload}
\end{figure}

\begin{figure}[H]
    \centering
    \includegraphics[width=1\linewidth]{images/admin_demo_02_import_results.png}
    \caption{Import results showing success notifications with session count. Administrators receive immediate feedback on upload success or errors.}
    \label{fig:admin-demo-results}
\end{figure}

\begin{figure}[H]
    \centering
    \includegraphics[width=1\linewidth]{images/admin_demo_03_import_history.png}
    \caption{Import history dashboard showing latest files uploaded for each category, with timestamps and file sizes.}
    \label{fig:admin-demo-history}
\end{figure}

\subsection{Authentication Interface Examples}

\begin{figure}[H]
    \centering
    \includegraphics[width=1\linewidth]{images/auth_demo_01_login.png}
    \caption{Login page with email and password fields. Users can sign in with existing credentials or navigate to registration.}
    \label{fig:auth-demo-login}
\end{figure}

\begin{figure}[H]
    \centering
    \includegraphics[width=1\linewidth]{images/auth_demo_02_signup.png}
    \caption{Registration page allowing new users to create accounts with email, password, and optional name fields.}
    \label{fig:auth-demo-signup}
\end{figure}

\section{System Performance and Metrics}

The system demonstrates strong performance characteristics:

\begin{center}
\begin{tabular}{|p{4cm}|p{2cm}|p{8cm}|}
\hline
\textbf{Metric} & \textbf{Value} & \textbf{Notes} \\
\hline
\textbf{Intent Classification} & 50-150ms & Deterministic classification typically \textless 50ms; Groq-based \textit{up to} 150ms. \\
\hline
\textbf{SQL Generation} & 200-500ms & Includes entity validation and query planning. \\
\hline
\textbf{Database Query} & 10-100ms & Depends on result set size and index effectiveness. \\
\hline
\textbf{Excel Import (50 sessions)} & 2-5s & Includes parsing, validation, and database insertion. \\
\hline
\textbf{Response Formatting} & 1-50ms & Heuristic rules typically fast; LLM polishing up to 50ms. \\
\hline
\textbf{Total Chat Response} & 300-800ms & End-to-end from submission to display. \\
\hline
\textbf{Deterministic Path \%} & 70-80\% & Most queries resolved without LLM calls. \\
\hline
\textbf{LLM Calls Avoided} & \textcolor{green}{~Significant} & Reduces latency and cost compared to naive approaches. \\
\hline
\end{tabular}
\end{center}

\section{Testing and Quality Assurance}

The backend includes comprehensive test coverage across multiple layers:

\begin{itemize}
    \item \textbf{Unit Tests}: Test individual services in isolation (intent\_router, sql\_agent, excel\_parser, etc.).
    \item \textbf{Integration Tests}: Test API endpoints with real database connections.
    \item \textbf{Import Tests}: Validate Excel parsing and database loading with sample workbooks.
    \item \textbf{Smoke Tests}: Quick validation of critical workflows.
    \item \textbf{Manual Testing}: Exploratory testing of UI interactions and edge cases.
\end{itemize}

Test files located in \texttt{backend/tests/}:
\begin{itemize}
    \item \texttt{test\_chat\_api.py}: Chat endpoint and routing logic.
    \item \texttt{test\_intent\_router.py}: Intent classification and entity extraction.
    \item \texttt{test\_sql\_agent.py}: SQL generation and query execution.
    \item \texttt{test\_admin\_import\_service.py}: Excel import workflow.
    \item \texttt{test\_excel\_parser.py}: Excel parsing and data transformation.
    \item \texttt{test\_auth\_service.py}: Authentication and authorization.
    \item \texttt{test\_groq\_service.py}: Language model integration.
    \item \texttt{test\_university\_info\_service.py}: Institutional information retrieval.
\end{itemize}

\section{Strengths and Technical Achievements}

The implementation demonstrates several significant strengths:

\begin{center}
\begin{tabular}{|p{3.5cm}|p{10cm}|}
\hline
\textbf{Strength} & \textbf{Description} \\
\hline

\textbf{Hybrid Architecture} & 
Combines deterministic logic (70-80\% of queries) with selective AI assistance, balancing speed, cost, and accuracy. \\
\hline

\textbf{Modular Services} & 
Clear service boundaries (IntentRouter, SQLAgent, GroqService, etc.) enable independent testing and evolution. \\
\hline

\textbf{Grounded Responses} & 
All timetable answers derived from database records; institutional information from official websites. No speculative generation. \\
\hline

\textbf{Safe Database Interaction} & 
Parameterized queries, SELECT-only enforcement, and validation prevent SQL injection and data corruption. \\
\hline

\textbf{Robust ETL Pipeline} & 
Automatic workbook category detection, comprehensive validation, transactional loading, and detailed error reporting. \\
\hline

\textbf{Role-Based Access Control} & 
Student, professor, and admin roles with fine-grained permission enforcement at the API level. \\
\hline

\textbf{Modern UI/UX} & 
React-based interface with real-time updates, conversation history, suggestion system, and responsive design. \\
\hline

\textbf{Multilingual Support} & 
Handles French, Arabic, and Tunisian Arabic (Arabizi) with typo tolerance and mixed-language queries. \\
\hline

\textbf{Real-Time Processing} & 
Fast intent classification and response generation enable natural conversational flow. \\
\hline

\textbf{Comprehensive Security} & 
Password hashing (PBKDF2-HMAC-SHA256), JWT tokens, session management, and structured access control. \\
\hline

\end{tabular}
\end{center}

\section{Limitations and Future Improvements}

While the system achieves its primary objectives, several limitations and opportunities for enhancement exist:

\begin{center}
\begin{tabular}{|p{3.5cm}|p{5.5cm}|p{4cm}|}
\hline
\textbf{Limitation} & \textbf{Impact} & \textbf{Mitigation Strategy} \\
\hline

\textbf{Dependency on Data Quality} & 
System accuracy depends on import data quality. Typos or inconsistencies propagate. & 
Advanced Excel validation rules and admin review workflows. \\
\hline

\textbf{Frontend Test Coverage} & 
Limited automated tests for UI components and interactions. & 
Expand Vitest coverage for React components. \\
\hline

\textbf{External Service Dependency} & 
Groq API availability affects system performance. Network latency on LLM calls. & 
Local fallback models; caching; rate limiting. \\
\hline

\textbf{Documentation Drift} & 
Documentation may lag behind active implementation. & 
Automated documentation generation from code. \\
\hline

\textbf{No Mobile App} & 
System is web-based only; no native mobile application. & 
Consider React Native or Flutter for mobile. \\
\hline

\textbf{Limited Personalization} & 
System treats all users the same; no user preferences or notification customization. & 
Add user preference storage and notification service. \\
\hline

\end{tabular}
\end{center}

\subsection{Recommended Future Work}

\begin{enumerate}
    \item \textbf{Observability}: Implement structured logging, distributed tracing, and metrics collection (Prometheus/Grafana).
    
    \item \textbf{Caching Strategy}: Implement multi-level caching (Redis for session data, HTTP caching for API responses).
        
    \item \textbf{Voice Interaction}: Add speech-to-text and text-to-speech for hands-free interaction.
    
    \item \textbf{Notification Service}: Implement push notifications for schedule changes, deadlines, and announcements.
    
    \item \textbf{Advanced Analytics}: Track user query patterns and system performance metrics for continuous improvement.
    
    \item \textbf{Multi-Language UI}: Localize interface text for French, Arabic, and English.
    
    \item \textbf{Integration with University Systems}: Connect with existing student management systems (APOGEE, STS, etc.).
\end{enumerate}

\section{Conclusion}

The implementation of Time Guide AI demonstrates a sophisticated, production-ready academic information system that integrates web development, database design, artificial intelligence, and user experience design into a coherent whole.

The system successfully translates theoretical foundations into practical outcomes:

\begin{itemize}
    \item \textbf{Conversational interfaces} enable non-technical users to access timetable information intuitively.
    \item \textbf{Hybrid AI architecture} balances flexibility with reliability through strategic use of deterministic logic and selective model assistance.
    \item \textbf{Modular services} enable maintainability and independent evolution of system components.
    \item \textbf{Grounded answering} ensures accuracy and trust through database records and official institutional sources.
    \item \textbf{Robust ETL pipeline} enables administrators to keep the system's knowledge base synchronized with operational data.
    \item \textbf{Security controls} protect user accounts and prevent unauthorized access to sensitive data.
\end{itemize}

The captured demonstrations show that the system effectively addresses the stated problem: transforming static Excel timetables into an intelligent, conversational, and reliable assistant that reduces friction for students, teachers, and administrators seeking university information.

The combination of technical rigor, thoughtful design, and practical usability makes Time Guide AI a strong foundation for institutional information systems in academic settings.
